\documentclass[10pt,a4paper]{amsart}
\usepackage[margin=19mm]{geometry}
\usepackage{amsmath,amssymb,mathtools}
\usepackage[scr=rsfs,cal=euler]{mathalfa}
\usepackage{xfrac}
\usepackage[unicode,hidelinks,bookmarks=false]{hyperref}
\newcommand{\ie}{i.e.\ }
\newcommand{\cf}{cf.\ }
\newcommand{\resp}{respectively\ }
\newcommand{\Subsection}{\subsection}
\providecommand{\nobreakdash}{\nobreak\hbox{-}}
\setlength{\emergencystretch}{2em}
\multlinegap0pt
\usepackage{amssymb}
\usepackage{mathtools}
\usepackage[shortlabels]{enumitem}
\usepackage{xcolor}
\newtheorem{proposition}{Proposition}
\newcommand{\id}{\operatorname{id}}
\title{Cocommutative Hopf Dialgebras and Rack Combinatorics}
\author{Jos\'e Gregorio Rodr\'iguez-Nieto, Olga Patricia Salazar-D\'iaz, Andr\'es Sarrazola-Alzate, Ra\'ul Vel\'asquez}
\date{Source: arXiv:2605.12749v1 (CC BY 4.0)}
\begin{document}
\maketitle

\noindent\emph{Source and licence.} Cocommutative Hopf Dialgebras and Rack Combinatorics. Version arXiv:2605.12749v1 (CC BY 4.0), distributed by arXiv under the Creative Commons Attribution 4.0 International licence, http://creativecommons.org/licenses/by/4.0/, as recorded in the arXiv licence metadata for this exact version and shown on the abstract page https://arxiv.org/abs/2605.12749v1. Full licence notice: \texttt{licenses/arxiv-2605.12749-CC-BY-4.0.txt}.\par\smallskip

\noindent\textit{Editorial scope.} Counted unit: the article's proposition prop:digroup-algebra-hopf-dialgebra (source0.tex lines 2182--2207). The article's complete original proof of it is delivered with it (source0.tex lines 2209--2364, one \texttt{proof} environment of 156 lines). No mathematics is written, repaired or re-derived: the statement and the proof are the article's own lines, in the article's own notation, and no retained line is substituted or re-flowed.  No further source-local context is needed: every cross-reference of every retained line resolves inside the two delivered spans.  Nothing was omitted from any retained span, so the package carries no omission record.  No retained line contains a citation, so the wrapper carries no bibliography and no bibliography entry.  No retained line contains a figure environment, an image-inclusion command or drawing code, so no image is delivered and the package declares no asset dependency.  Editorial formatting only, in addition: an AMS \texttt{amsart} preamble that restates the article's own declarations and macros used by the retained lines, namely the environments \texttt{proposition} on separate counters, together with the standard packages those lines need.  The retained environments print in this document as Proposition 1 in order of appearance, whereas the article prints the same items with the numbering of its own full document.  The complete native source of the article as distributed in the arXiv e-print for this exact version is bundled unchanged as \texttt{sources/200445/source0.tex}.
	\begin{proposition}\label{prop:digroup-algebra-hopf-dialgebra}
		Let \(D\) be a digroup with products \(\vdash,\dashv\), distinguished
		bar-unit \(\xi\), and inverse map \(x\mapsto x^{-1}\). Let
		\[
		K[D]:=\bigoplus_{x\in D}Kx
		\]
		be the free \(K\)-vector space with basis \(D\). Extend the two products
		of \(D\) bilinearly to maps
		\[
		\vdash,\dashv:K[D]\otimes K[D]\longrightarrow K[D].
		\]
		For \(x\in D\), define
		\[
		\Delta(x)=x\otimes x,
		\qquad
		\varepsilon(x)=1,
		\qquad
		S(x)=x^{-1},
		\]
		and extend \(\Delta\), \(\varepsilon\), and \(S\) linearly to \(K[D]\).
		Then
		\[
		\bigl(K[D],\Delta,\varepsilon,\xi,\vdash,\dashv,S\bigr)
		\]
		is a cocommutative Hopf dialgebra.
	\end{proposition}

	\begin{proof}
		We verify the defining properties on basis elements. The corresponding
		identities on all of \(K[D]\) then follow by multilinearity.
		
		First, since \(D\) is a digroup, both \((D,\vdash)\) and \((D,\dashv)\)
		are semigroups, and the mixed dialgebra identities
		\[
		x\vdash (y\dashv z)=(x\vdash y)\dashv z,
		\]
		\[
		x\dashv (y\dashv z)=x\dashv (y\vdash z),
		\]
		and
		\[
		(x\vdash y)\vdash z=(x\dashv y)\vdash z
		\]
		hold for all \(x,y,z\in D\). Extending the products bilinearly, these same
		identities hold on \(K[D]\). Hence \(K[D]\) is a dialgebra.
		
		The distinguished bar-unit \(\xi\in D\) satisfies
		\[
		\xi\vdash x=x=x\dashv \xi
		\qquad
		(x\in D).
		\]
		Moreover, the balanced identity
		\[
		x\vdash \xi=\xi\dashv x
		\qquad
		(x\in D)
		\]
		follows from the inverse identities. Indeed, since
		\[
		x^{-1}\dashv x=\xi,
		\]
		we have
		\[
		x\vdash \xi
		=
		x\vdash (x^{-1}\dashv x)
		=
		(x\vdash x^{-1})\dashv x
		=
		\xi\dashv x.
		\]
		Thus \(K[D]\) is a balanced bar-unital dialgebra.
		
		We now check the coalgebra structure. For \(x\in D\),
		\[
		(\Delta\otimes \id)\Delta(x)
		=
		(\Delta\otimes \id)(x\otimes x)
		=
		x\otimes x\otimes x,
		\]
		and similarly
		\[
		(\id\otimes \Delta)\Delta(x)
		=
		x\otimes x\otimes x.
		\]
		Hence \(\Delta\) is coassociative. Also,
		\[
		(\varepsilon\otimes \id)\Delta(x)=\varepsilon(x)x=x,
		\qquad
		(\id\otimes \varepsilon)\Delta(x)=x\varepsilon(x)=x.
		\]
		Therefore \((K[D],\Delta,\varepsilon)\) is a coalgebra. It is
		cocommutative because
		\[
		\tau\Delta(x)=\tau(x\otimes x)=x\otimes x=\Delta(x).
		\]
		Furthermore,
		\[
		\Delta(\xi)=\xi\otimes\xi,
		\qquad
		\varepsilon(\xi)=1,
		\]
		so \(\xi\) is group-like.
		
		Next we show that the two products are coalgebra morphisms. Let
		\(\star\in\{\vdash,\dashv\}\). For \(x,y\in D\), one has
		\[
		\Delta(x\star y)
		=
		(x\star y)\otimes(x\star y).
		\]
		On the other hand,
		\begin{align*}
			(\star\otimes \star)
			\circ(\id\otimes\tau\otimes\id)
			\bigl(\Delta(x)\otimes\Delta(y)\bigr)
			&=
			(\star\otimes \star)
			\circ(\id\otimes\tau\otimes\id)
			\bigl((x\otimes x)\otimes(y\otimes y)\bigr)\\
			&=
			(\star\otimes \star)(x\otimes y\otimes x\otimes y)\\
			&=
			(x\star y)\otimes(x\star y).
		\end{align*}
		Thus \(\Delta\circ\star\) satisfies the required compatibility. Similarly,
		\[
		\varepsilon(x\star y)=1=\varepsilon(x)\varepsilon(y).
		\]
		Hence both \(\vdash\) and \(\dashv\) are coalgebra morphisms.
		
		Now consider the antipode map. Since \(S(x)=x^{-1}\), we have
		\[
		\Delta(S(x))
		=
		\Delta(x^{-1})
		=
		x^{-1}\otimes x^{-1}
		=
		(S\otimes S)(x\otimes x)
		=
		(S\otimes S)\Delta(x),
		\]
		and
		\[
		\varepsilon(S(x))=\varepsilon(x^{-1})=1=\varepsilon(x).
		\]
		Thus \(S\) is a coalgebra morphism.
		
		Finally, the inverse identities in \(D\) give, for every \(x\in D\),
		\[
		\sum_{(x)}x_{(1)}\vdash S(x_{(2)})
		=
		x\vdash S(x)
		=
		x\vdash x^{-1}
		=
		\xi
		=
		\varepsilon(x)\xi,
		\]
		and
		\[
		\sum_{(x)}S(x_{(1)})\dashv x_{(2)}
		=
		S(x)\dashv x
		=
		x^{-1}\dashv x
		=
		\xi
		=
		\varepsilon(x)\xi.
		\]
		By linearity, both antipode identities hold for every element of \(K[D]\).
		Therefore
		\[
		\bigl(K[D],\Delta,\varepsilon,\xi,\vdash,\dashv,S\bigr)
		\]
		is a cocommutative Hopf dialgebra.
	\end{proof}

\end{document}
